\chapter{Spectral--Operator Compatibility as a Model-Adequacy Problem}
\label{ch:proof-architecture}

\section{The central question}

A reduced Hamiltonian may be judged by more than one criterion.  A spectral reduction
asks whether selected eigenvalues or derived band quantities are reproduced within
declared tolerances.  An operator reduction asks whether the reduced Hamiltonian
reproduces matrix information in a common aligned representation.  These requirements
are not equivalent.  Two models may have nearly identical selected eigenvalues while
differing in onsite terms, orbital couplings, hopping blocks, or their response to a
later perturbation.

The central model-adequacy question is therefore not whether two independently optimized
models are individually successful.  It is whether one candidate Hamiltonian satisfies
both requirements.

Let $J$ be a declared finite ordered inventory of model classes.  For each $j\in J$,
let $\mathfrak M_j$ be the corresponding candidate class and let $\theta$ denote a
candidate parameter in that class.  Write
\[
  H_j(\theta)\in\mathfrak M_j
\]
for the candidate Hamiltonian.  A spectral loss $\mathcal L_{E,j}(\theta)$ and an
aligned operator loss $\mathcal L_{H,j}(C,\theta)$ are evaluated against one declared
reference problem.  The map $C$ identifies candidate coordinates with reference
coordinates and must belong to a specified admissible alignment family
$\mathcal U_j$.

The two admissible sets are
\[
  \mathfrak A_{E,j}(\tau_E)
  =\left\{
      H_j(\theta)\in\mathfrak M_j:
      \mathcal L_{E,j}(\theta)\leq\tau_E
    \right\}
\]
and
\[
  \mathfrak A_{H,j}(\tau_H)
  =\left\{
      H_j(\theta)\in\mathfrak M_j:
      \inf_{C\in\mathcal U_j}
      \mathcal L_{H,j}(C,\theta)\leq\tau_H
    \right\}.
\]
When the alignment optimization is known to attain its infimum, the infimum may be
written as a minimum.  Attainment is a mathematical premise, not a consequence of
running an optimizer.

Exact compatibility within the model class is the proposition
\[
  \mathfrak A_{E,j}(\tau_E)
  \cap
  \mathfrak A_{H,j}(\tau_H)
  \neq\varnothing.
\]
When the alignment infimum is attained, this is equivalently the existence of one
parameter $\theta$ and one admissible alignment $C\in\mathcal U_j$ such that
\[
  \mathcal L_{E,j}(\theta)\leq\tau_E,
  \qquad
  \mathcal L_{H,j}(C,\theta)\leq\tau_H.
\]
The quantifiers are coupled.  A spectrally successful parameter, an operator-successful
parameter, and a favorable alignment chosen for some third parameter do not assemble
into a common witness.

This is the problem stated by the ICMSEP extended abstract that motivates the present
formal synthesis: determine whether spectral and operator reductions admit a common
Hamiltonian in each prescribed model class, quantify their separation when they do not,
and identify the first tested class that supports a common endpoint.  The present
companion develops exact finite mathematics supporting that question.  It does not
claim that the proposed bulk-silicon calculation has been executed or scientifically
validated.

\section{Separation and tolerance-level compatibility}

When the two admissible sets do not intersect, a metric on common canonical model
coordinates can measure the obstruction.  If $d_{\mathrm{model}}$ is that declared
metric, define
\[
  \delta_j^*
  =\inf_{\substack{
       H_E\in\mathfrak A_{E,j}(\tau_E)\\
       H_H\in\mathfrak A_{H,j}(\tau_H)}}
    d_{\mathrm{model}}(H_E,H_H).
\]
A positive certified lower bound on $\delta_j^*$ proves that the two sets are disjoint.
The converse direction requires care.  In an arbitrary metric space, distance zero does
not by itself produce an intersection point.  Exact equivalence between zero distance
and nonempty intersection requires suitable closure, compactness, attainment, or other
explicit hypotheses.  The finite geometric results used later state the assumptions
under which this converse is valid.

An application may also introduce a map tolerance $\tau_{\mathrm{map}}$ and declare
$\delta_j^*\leq\tau_{\mathrm{map}}$ to be an accepted approximate compatibility
condition.  This is a third declared criterion.  It must not be silently identified
with literal set intersection.  The exact intersection proposition, the exact set
distance, and the tolerance-level decision remain distinct even when they are reported
together.

\section{The model hierarchy}

The ICMSEP problem is indexed by a hierarchy of candidate classes rather than by one
unconstrained matrix space.  A lower-order class may impose a fixed orbital basis,
short hopping range, or restricted symmetry parameterization; a larger class may add
neighbor shells or selected symmetry-allowed corrections.  The scientific decision is
the first tested class for which the declared compatibility predicate passes.

Nested model classes do not automatically make every computed loss monotone.  Enlarging
$\mathfrak M_j$ preserves existential feasibility only when the class inclusion, the
reference data, the losses, the tolerances, and the alignment-family embedding are
conformant.  Refitting, changing an optimizer basin, changing support, or changing the
alignment family can break naive monotonicity of reported minimizers.  The finite
first-passage logic developed in Chapter~\ref{ch:reduction-decision} therefore examines
every earlier declared class directly rather than assuming monotone numerical behavior.

The phrase ``smallest adequate model'' is consequently shorthand for a precise finite
statement: the selected class passes the joint predicate, and every earlier class in the
declared inventory fails it.  It is not a claim about all conceivable Hamiltonian
families.

\section{The alignment problem is a component of operator admissibility}

The operator loss cannot be evaluated until the reference and candidate coordinates
have been identified.  For momentum-indexed retained operators, the aligned difference
used in this companion is
\[
  \Delta_{C,\theta}(k)
  =H_{\mathrm{ref}}(k)
   -C(k)^{-1}H_j(\theta;k)C(k).
\]
A weighted Frobenius operator loss may then be formed from momentum-space differences or
from corresponding real-space blocks.  The exact choice of weights, blocks, and
normalization belongs to the definition of $\mathcal L_{H,j}$.

The admissible alignment family $\mathcal U_j$ is part of that definition.  Identity
alignment, one global unitary, arbitrary pointwise unitaries, and symmetry-constrained
or finite-depth families produce different operator-admissible sets.  Enlarging the
family can reduce the optimized operator loss, but it can also admit transformations
that alter Fourier support or violate an application's structural requirements.

Thus the family-relative question
\[
  \exists C\in\mathcal U_j,
  \qquad
  \mathcal L_{H,j}(C,\theta)\leq\tau_H
\]
is a fiber of the full compatibility problem at one candidate $\theta$.  It is not a
replacement for the joint quantification over candidate Hamiltonians.  Chapters
\ref{ch:finite-operator-foundations} and \ref{ch:bloch-alignment-locality} develop this
operator branch because it is necessary for defining $\mathfrak A_{H,j}$ correctly.

\section{Fourier locality is a structural diagnostic}

For a momentum-dependent retained Hamiltonian, finite Fourier transformation produces
real-space blocks.  Hopping range, blockwise residuals, and truncation error are
therefore properties of the operator representation.  A unitary common to all momentum
points passes through the Fourier sum, while a momentum-dependent unitary can
redistribute Fourier weight among modes.  Consequently, a pointwise alignment that
improves the operator loss may change locality diagnostics.

The Fourier results in this companion establish three bounded facts.  Fixed-unitary
conjugation commutes with the finite Fourier operations used here.  An exact two-point
rank-two construction disproves universal commutation between pointwise conjugation and
support truncation.  Exact finite-tail identities and inequalities then control the
deliberate-truncation channel for a declared support.

These results clarify how $\mathcal L_{H,j}$ and a locality or support requirement may
interact.  They do not define the spectral admissible set, select a material model, or
replace the full spectral--operator intersection problem.

\section{Four interacting mathematical layers}

The proof architecture can now be stated in its proper order.

\subsection{Candidate and criterion layer}

A declared model class, reference problem, spectral loss, operator loss, tolerances, and
alignment family determine the two admissible sets.  Their intersection is the primary
compatibility proposition.  Their metric separation quantifies an obstruction under
additional geometric hypotheses.

\subsection{Representation and alignment layer}

Finite frames represent retained subspaces.  Projectors are invariant under internal
right-unitary changes, compressed operators transform covariantly, and reconstructed
ambient operators are invariant.  Between two models, an admissible alignment pulls the
candidate into reference coordinates before an operator residual is formed.

\subsection{Fourier and locality layer}

Finite Fourier transformation converts the aligned operator field into real-space
blocks.  Fixed and pointwise coordinate changes have different support consequences.
For a reference field, an interpolant, and its deliberately truncated approximation,
the interpolation, truncation, and total residuals form an exact additive chain at the
vector level.  Their norms satisfy triangle bounds rather than a scalar additive
identity.  Only the deliberate-truncation leg is identified with the checked omitted
Fourier tail.

\subsection{Geometry, certificate, and decision layer}

Exact admissible-set geometry proves intersection or separation for the families covered
by its hypotheses.  Certificate checkers connect finite data to exact propositions.
Three-way interval decisions preserve unresolved threshold overlap, and finite
first-passage logic selects the earliest class in one declared inventory.

\section{Organization}

Chapter~\ref{ch:finite-operator-foundations} proves the representation-neutral frame,
projector, compression, and reconstruction identities needed by the aligned operator
criterion.  Chapter~\ref{ch:bloch-alignment-locality} adds momentum dependence, aligned
operator differences, alignment families, fixed and pointwise gauge actions, the
three-channel residual calculus, the rank-two noncommutation witness, and exact
Fourier-tail control.
Chapter~\ref{ch:reduction-decision} returns to the full spectral--operator compatibility
problem and develops admissible-set geometry, separation, certificate semantics,
interval dispositions, and selection over a finite model inventory.
Chapter~\ref{ch:mechanization-status} states what has and has not yet been assembled into
an end-to-end formal adequacy theorem.

The appendices provide declaration-level traceability, immutable provenance, audit
boundaries, notation bridges, and a source- and status-qualified crosswalk of the
error-channel development across the seven-paper program.  They separate repository
mechanics from the mathematical narrative without changing the authority of the
underlying Lean statements.

\section{Difference convention and scope}

Throughout the companion, the aligned difference is reference minus pulled-back
candidate:
\[
  \Delta_{\mathrm{ref},\mathrm{cand}}(k)
  =H_{\mathrm{ref}}(k)-C(k)^{-1}H_{\mathrm{cand}}(k)C(k).
\]
Reversing the operands negates the difference.  Norm-based losses are insensitive to
this sign, but signed perturbations, defect potentials, and response coefficients are
not.

All results concern exact finite objects under explicit hypotheses.  They do not by
themselves establish that a runtime frame is exactly orthonormal, that floating-point
intervals contain exact quantities, that a sampled spectral gap persists between sample
points, or that a candidate is physically adequate for silicon.  In particular, the
companion supplies formal components of the ICMSEP decision architecture; it does not
claim completion of the proposed first-principles calculation, a retained certificate,
or a material-validation result.
